% =====================================================================
% 04-Topology.tex
% Geometric Hallucination Taxonomy — Mathematical Spine
% Borealis S. Hedling — 2026
% =====================================================================

\section*{04. Topological Operator}

This module formalizes metric rank, degeneracy, collapse surfaces, and
instability criteria used to define Topological Instability Hallucinations
(TIH) in the Geometric Hallucination Taxonomy (GHT). All definitions rely
on the geometric foundations established in 00-Foundations.tex.

% ---------------------------------------------------------------------
\subsection*{4.1 Metric Tensor and Rank}

Let $g$ be the Riemannian metric on $M$. At each point $p \in M$, the
metric defines a symmetric, positive–definite bilinear form
\[
g_p : T_p M \times T_p M \to \mathbb{R}.
\]

In local coordinates, $g_p$ is represented by the matrix $(g_{ij}(p))$.

The rank of the metric at $p$ is
\[
\tau(p) = \operatorname{rank}(g_p).
\]

For a nondegenerate Riemannian metric, $\tau(p) = n$ everywhere.  
Topological instability arises when $\tau(p)$ drops below full rank.

% ---------------------------------------------------------------------
\subsection*{4.2 Degenerate Regions}

A point $p$ is degenerate if
\[
\det(g_p) = 0.
\]

The set of degenerate points is the collapse surface
\[
\Sigma = \{ p \in M : \det(g_p) = 0 \}.
\]

If $\Sigma$ is nonempty, the manifold contains regions where the metric
fails to define a valid inner product.

% ---------------------------------------------------------------------
\subsection*{4.3 Collapse Surfaces}

A collapse surface $\Sigma$ may be:
\begin{itemize}
\item a hypersurface (codimension 1),
\item a lower–dimensional submanifold,
\item a fractal or singular set (in pathological cases).
\end{itemize}

Crossing $\Sigma$ corresponds to a topological phase transition in the
latent geometry.

% ---------------------------------------------------------------------
\subsection*{4.4 Rank Instability}

Rank instability occurs when
\[
\tau(p) < n.
\]

A latent trajectory $x(t)$ exhibits rank instability at time $t$ if
\[
\tau(x(t)) < n.
\]

Oscillatory rank instability occurs when
\[
\tau(x(t+1)) \neq \tau(x(t)).
\]

Collapse instability occurs when
\[
\tau(x(t)) \to 0,
\]
indicating total metric degeneration.

% ---------------------------------------------------------------------
\subsection*{4.5 Collapse Dynamics Along Trajectories}

Given a latent trajectory $x(t)$, define the collapse indicator
\[
C(t) = \det(g_{x(t)}).
\]

Collapse occurs when
\[
C(t) \to 0.
\]

Collapse rate is defined by
\[
\Delta_C(t) = C(t+1) - C(t).
\]

A negative collapse rate indicates movement toward degeneracy.

% ---------------------------------------------------------------------
\subsection*{4.6 Topological Stability Region}

Define the topological stability region
\[
\Omega_{\tau}
= \{ p \in M : \tau(p) = n \}.
\]

A latent trajectory is topologically stable if $x(t) \in \Omega_{\tau}$
for all $t$.

% ---------------------------------------------------------------------
\subsection*{4.7 Topological Instability Hallucinations (TIH)}

A latent trajectory exhibits TIH if there exists a sequence of times $t$
such that
\[
\tau(x(t)) < n
\quad \text{for infinitely many } t.
\]

Equivalently, TIH occurs when metric degeneracy prevents the latent
trajectory from remaining within a nondegenerate geometric region.

TIH corresponds to structural collapse of the latent manifold.

% ---------------------------------------------------------------------
\subsection*{4.8 Discrete Approximation}

In numerical settings, the metric matrix is approximated by
\[
g_{ij}(x_t)
= \langle e_i, e_j \rangle_{g_{x_t}},
\]
where $\{e_i\}$ is a coordinate basis.

Rank is computed via singular value decomposition (SVD):
\[
\tau(x_t) = \#\{ \sigma_i > \epsilon \},
\]
where $\epsilon$ is a numerical tolerance.

The Python toolkit implements this approximation and validates it in
\texttt{/tests/test\_topology.py}.

% ---------------------------------------------------------------------
\subsection*{4.9 Collapse Surfaces and Hallucination Onset}

Collapse surfaces often correspond to hallucination onset regions.  
When the metric degenerates, distances, angles, and curvature become
ill-defined, causing latent representations to collapse into degenerate
subspaces.

This geometric phenomenon motivates the TIH classification.

% =====================================================================
% End of 04-Topology.tex
% =====================================================================
